\IfFileExists{stacks-project.cls}{%
\documentclass{stacks-project}
}{%
\documentclass{amsart}
}
\usepackage{amssymb}

% For dealing with references we use the comment environment
\usepackage{verbatim}
\usepackage{fancyvrb}
\newenvironment{reference}{\comment}{\endcomment}
%\newenvironment{reference}{}{}
\newenvironment{slogan}{\comment}{\endcomment}
\newenvironment{history}{\comment}{\endcomment}

% For source figures retained as assets
\usepackage{graphicx}

% For commutative diagrams we use Xy-pic
\usepackage[all]{xy}

% We use 2cell for 2-commutative diagrams.
\xyoption{2cell}
\UseAllTwocells

% We use multicol for the list of chapters between chapters
\usepackage{multicol}

% This is generally recommended for better output
\usepackage{lmodern}
\usepackage[T1]{fontenc}
\usepackage{textalpha}

% For cross-file-references
% We prefer the snapshot tmp/xref/#2.aux,
% which latexmkrc replaces atomically, so
% that documents compiled simultaneously
% never read a half-written .aux file.
% The last argument keeps links to #2.pdf.
\usepackage{xr-hyper}
\newcommand{\maybeexternaldocument}[2]{%
  \IfFileExists{#2.aux}{%
    \externaldocument[#1]{#2}%
  }{%
    \PackageWarning{stack}{Missing #2.aux; external references with prefix #1 unavailable}%
  }%
}
\let\externaldocumentorig\externaldocument
\renewcommand{\externaldocument}[2][]{%
  \IfFileExists{tmp/xref/#2.aux}{%
    \externaldocumentorig[#1]{tmp/xref/#2}[#2.pdf]%
  }{%
  \IfFileExists{#2.aux}{%
    \externaldocumentorig[#1]{#2}%
  }{%
    \PackageWarning{stack}{Missing #2.aux; external references with prefix #1 unavailable}%
  }%
  }%
}

% Package for hypertext links:
\usepackage{hyperref}

% For any local file, say "hello.tex" you want to link to please
% use \externaldocument[hello-]{hello}
\externaldocument[introduction-]{introduction}
\externaldocument[conventions-]{conventions}
\externaldocument[sets-]{sets}
\externaldocument[categories-]{categories}
\externaldocument[topology-]{topology}
\externaldocument[sheaves-]{sheaves}
\externaldocument[sites-]{sites}
\externaldocument[stacks-]{stacks}
\externaldocument[fields-]{fields}
\externaldocument[algebra-]{algebra}
\externaldocument[brauer-]{brauer}
\externaldocument[homology-]{homology}
\externaldocument[derived-]{derived}
\externaldocument[simplicial-]{simplicial}
\externaldocument[more-algebra-]{more-algebra}
\externaldocument[smoothing-]{smoothing}
\externaldocument[modules-]{modules}
\externaldocument[sites-modules-]{sites-modules}
\externaldocument[injectives-]{injectives}
\externaldocument[cohomology-]{cohomology}
\externaldocument[sites-cohomology-]{sites-cohomology}
\externaldocument[dga-]{dga}
\externaldocument[dpa-]{dpa}
\externaldocument[sdga-]{sdga}
\externaldocument[hypercovering-]{hypercovering}
\externaldocument[schemes-]{schemes}
\externaldocument[constructions-]{constructions}
\externaldocument[properties-]{properties}
\externaldocument[morphisms-]{morphisms}
\externaldocument[coherent-]{coherent}
\externaldocument[divisors-]{divisors}
\externaldocument[limits-]{limits}
\externaldocument[varieties-]{varieties}
\externaldocument[topologies-]{topologies}
\externaldocument[descent-]{descent}
\externaldocument[perfect-]{perfect}
\externaldocument[more-morphisms-]{more-morphisms}
\externaldocument[flat-]{flat}
\externaldocument[groupoids-]{groupoids}
\externaldocument[more-groupoids-]{more-groupoids}
\externaldocument[etale-]{etale}
\externaldocument[chow-]{chow}
\externaldocument[intersection-]{intersection}
\externaldocument[pic-]{pic}
\externaldocument[weil-]{weil}
\externaldocument[adequate-]{adequate}
\externaldocument[dualizing-]{dualizing}
\externaldocument[duality-]{duality}
\externaldocument[discriminant-]{discriminant}
\externaldocument[derham-]{derham}
\externaldocument[local-cohomology-]{local-cohomology}
\externaldocument[algebraization-]{algebraization}
\externaldocument[curves-]{curves}
\externaldocument[resolve-]{resolve}
\externaldocument[models-]{models}
\externaldocument[functors-]{functors}
\externaldocument[equiv-]{equiv}
\externaldocument[pione-]{pione}
\externaldocument[etale-cohomology-]{etale-cohomology}
\externaldocument[proetale-]{proetale}
\externaldocument[relative-cycles-]{relative-cycles}
\externaldocument[more-etale-]{more-etale}
\externaldocument[trace-]{trace}
\externaldocument[crystalline-]{crystalline}
\externaldocument[spaces-]{spaces}
\externaldocument[spaces-properties-]{spaces-properties}
\externaldocument[spaces-morphisms-]{spaces-morphisms}
\externaldocument[decent-spaces-]{decent-spaces}
\externaldocument[spaces-cohomology-]{spaces-cohomology}
\externaldocument[spaces-limits-]{spaces-limits}
\externaldocument[spaces-divisors-]{spaces-divisors}
\externaldocument[spaces-over-fields-]{spaces-over-fields}
\externaldocument[spaces-topologies-]{spaces-topologies}
\externaldocument[spaces-descent-]{spaces-descent}
\externaldocument[spaces-perfect-]{spaces-perfect}
\externaldocument[spaces-more-morphisms-]{spaces-more-morphisms}
\externaldocument[spaces-flat-]{spaces-flat}
\externaldocument[spaces-groupoids-]{spaces-groupoids}
\externaldocument[spaces-more-groupoids-]{spaces-more-groupoids}
\externaldocument[bootstrap-]{bootstrap}
\externaldocument[spaces-pushouts-]{spaces-pushouts}
\externaldocument[spaces-chow-]{spaces-chow}
\externaldocument[groupoids-quotients-]{groupoids-quotients}
\externaldocument[spaces-more-cohomology-]{spaces-more-cohomology}
\externaldocument[spaces-simplicial-]{spaces-simplicial}
\externaldocument[spaces-duality-]{spaces-duality}
\externaldocument[formal-spaces-]{formal-spaces}
\externaldocument[restricted-]{restricted}
\externaldocument[spaces-resolve-]{spaces-resolve}
\externaldocument[formal-defos-]{formal-defos}
\externaldocument[defos-]{defos}
\externaldocument[cotangent-]{cotangent}
\externaldocument[examples-defos-]{examples-defos}
\externaldocument[algebraic-]{algebraic}
\externaldocument[examples-stacks-]{examples-stacks}
\externaldocument[stacks-sheaves-]{stacks-sheaves}
\externaldocument[criteria-]{criteria}
\externaldocument[artin-]{artin}
\externaldocument[quot-]{quot}
\externaldocument[stacks-properties-]{stacks-properties}
\externaldocument[stacks-morphisms-]{stacks-morphisms}
\externaldocument[stacks-limits-]{stacks-limits}
\externaldocument[stacks-cohomology-]{stacks-cohomology}
\externaldocument[stacks-perfect-]{stacks-perfect}
\externaldocument[stacks-introduction-]{stacks-introduction}
\externaldocument[stacks-more-morphisms-]{stacks-more-morphisms}
\externaldocument[stacks-geometry-]{stacks-geometry}
\externaldocument[moduli-]{moduli}
\externaldocument[moduli-curves-]{moduli-curves}
\externaldocument[examples-]{examples}
\externaldocument[exercises-]{exercises}
\externaldocument[guide-]{guide}
\externaldocument[desirables-]{desirables}
\externaldocument[coding-]{coding}
\externaldocument[obsolete-]{obsolete}
\externaldocument[fdl-]{fdl}
\externaldocument[index-]{index}

% Heap Project documents
\externaldocument[k3-]{k3}
\externaldocument[torelli-]{torelli}
\externaldocument[linear-algebra-]{linear-algebra}
\externaldocument[groups-]{groups}
\externaldocument[complex-analysis-]{complex-analysis}
\externaldocument[functional-analysis-]{functional-analysis}
\externaldocument[categories-]{categories}
\externaldocument[commutative-algebra-]{commutative-algebra}
\externaldocument[ringed-spaces-]{ringed-spaces}
\externaldocument[homological-algebra-]{homological-algebra}
\externaldocument[lie-algebras-]{lie-algebras}
\externaldocument[representation-theory-]{representation-theory}
\externaldocument[smooth-manifolds-]{smooth-manifolds}
\externaldocument[differential-topology-]{differential-topology}
\externaldocument[algebraic-topology-]{algebraic-topology}
\externaldocument[differential-geometry-]{differential-geometry}
\externaldocument[algebraic-geometry-]{algebraic-geometry}
\externaldocument[symplectic-geometry-]{symplectic-geometry}
\externaldocument[complex-geometry-]{complex-geometry}
\externaldocument[hodge-theory-]{hodge-theory}
\externaldocument[%
variations-of-hodge-structures-]{%
variations-of-hodge-structures}
\externaldocument[%
homological-mirror-symmetry-]{%
homological-mirror-symmetry}
\maybeexternaldocument{deformations-}{deformations}
\externaldocument[vertex-algebras-]{vertex-algebras}
\externaldocument[springer-]{springer}
\externaldocument[infty-categories-]{infty-categories}
\externaldocument[seminars-differential-geometry-]{%
seminars-differential-geometry}
\externaldocument[seminars-complex-geometry-]{%
seminars-complex-geometry}
\externaldocument[seminars-algebraic-geometry-]{%
seminars-algebraic-geometry}
\externaldocument[seminars-others-]{%
seminars-others}
\externaldocument[%
geometric-prequantization-]{%
geometric-prequantization}
\externaldocument[surfaces-]{surfaces}
\externaldocument[birational-maps-commutative-algebra-]{birational-maps-commutative-algebra}
\externaldocument[cremona-transformations-]{cremona-transformations}
\externaldocument[derived-categories-of-sheaves-]{derived-categories-of-sheaves}
\externaldocument[geometric-invariant-theory-]{geometric-invariant-theory}
\externaldocument[geometric-stability-conditions-and-group-actions-]{geometric-stability-conditions-and-group-actions}
\externaldocument[moduli-spaces-of-sheaves-]{moduli-spaces-of-sheaves}
\externaldocument[bridgeland-stability-general-theory-]{bridgeland-stability-general-theory}
\externaldocument[hilbert-schemes-]{hilbert-schemes}
\externaldocument[noncommutative-abelian-surfaces-and-kummer-type-hyperkahler-varieties-]{noncommutative-abelian-surfaces-and-kummer-type-hyperkahler-varieties}
\externaldocument[mumford-tate-groups-in-hodge-theory-]{mumford-tate-groups-in-hodge-theory}
\externaldocument[hodge-birational-atoms-2026-]{hodge-birational-atoms-2026}
\externaldocument[xxii-egd-2026-]{%
xxii-egd-2026}
\externaldocument[pde-]{pde}
\externaldocument[probability-]{probability}

\externaldocument[linguistics-]{linguistics}
\externaldocument[genealogia-familiar-]{%
genealogia-familiar}

\externaldocument[%
16th-alga-meeting-2026-]{%
16th-alga-meeting-2026}

\externaldocument[%
iv-jornadas-lefschetz-2026-]{%
iv-jornadas-lefschetz-2026}

\externaldocument[reading-the-masters-]{%
reading-the-masters}

\externaldocument[real-analysis-]{%
real-analysis}

\externaldocument[optimal-transport-]{%
optimal-transport}

\externaldocument[history-]{history}

\externaldocument[seiberg-witten-]{seiberg-witten}

\externaldocument[history-philosophy-biology-]{history-philosophy-biology}

\externaldocument[chemistry-]{chemistry}

\externaldocument[molecular-cell-biology-]{molecular-cell-biology}

\externaldocument[microbiology-]{microbiology}

\externaldocument[protistology-]{protistology}

\externaldocument[mycology-]{mycology}

\externaldocument[botany-]{botany}

\externaldocument[zoology-]{zoology}

\externaldocument[animal-histology-]{animal-histology}

\externaldocument[animal-physiology-]{animal-physiology}

\externaldocument[systematics-biogeography-]{systematics-biogeography}

\externaldocument[ecology-natural-resources-]{ecology-natural-resources}

\externaldocument[earth-sciences-]{earth-sciences}

\externaldocument[biostatistics-]{biostatistics}

\externaldocument[mathematical-systems-biology-]{mathematical-systems-biology}

\externaldocument[cimpa-el-salvador-2024-]{cimpa-el-salvador-2024}

% Theorem environments.
%
\theoremstyle{plain}
\newtheorem{theorem}[subsection]{Theorem}
\newtheorem{proposition}[subsection]{Proposition}
\newtheorem{lemma}[subsection]{Lemma}
\newtheorem{corollary}[subsection]{Corollary}

\theoremstyle{definition}
\newtheorem{definition}[subsection]{Definition}
\newtheorem{example}[subsection]{Example}
\newtheorem{exercise}[subsection]{Exercise}
\newtheorem{situation}[subsection]{Situation}

\theoremstyle{remark}
\newtheorem{remark}[subsection]{Remark}
\newtheorem{remarks}[subsection]{Remarks}

\numberwithin{equation}{subsection}

% Macros
%
\def\lim{\mathop{\mathrm{lim}}\nolimits}
\def\colim{\mathop{\mathrm{colim}}\nolimits}
\def\Spec{\mathop{\mathrm{Spec}}}
\def\Hom{\mathop{\mathrm{Hom}}\nolimits}
\def\Ext{\mathop{\mathrm{Ext}}\nolimits}
\def\SheafHom{\mathop{\mathcal{H}\!\mathit{om}}\nolimits}
\def\SheafExt{\mathop{\mathcal{E}\!\mathit{xt}}\nolimits}
\def\Sch{\mathit{Sch}}
\def\Mor{\mathop{\mathrm{Mor}}\nolimits}
\def\Ob{\mathop{\mathrm{Ob}}\nolimits}
\def\Sh{\mathop{\mathit{Sh}}\nolimits}
\def\NL{\mathop{N\!L}\nolimits}
\def\CH{\mathop{\mathrm{CH}}\nolimits}
\def\proetale{{pro\text{-}\acute{e}tale}}
\def\etale{{\acute{e}tale}}
\def\QCoh{\mathit{QCoh}}
\def\Ker{\mathop{\mathrm{Ker}}}
\def\Im{\mathop{\mathrm{Im}}}
\def\Coker{\mathop{\mathrm{Coker}}}
\def\Coim{\mathop{\mathrm{Coim}}}

% Boxtimes
%
\DeclareMathSymbol{\boxtimes}{\mathbin}{AMSa}{"02}

%
% Macros for moduli stacks/spaces
%
\def\QCohstack{\mathcal{QC}\!\mathit{oh}}
\def\Cohstack{\mathcal{C}\!\mathit{oh}}
\def\Spacesstack{\mathcal{S}\!\mathit{paces}}
\def\Quotfunctor{\mathrm{Quot}}
\def\Hilbfunctor{\mathrm{Hilb}}
\def\Curvesstack{\mathcal{C}\!\mathit{urves}}
\def\Polarizedstack{\mathcal{P}\!\mathit{olarized}}
\def\Complexesstack{\mathcal{C}\!\mathit{omplexes}}
% \Pic is the operator that assigns to X its picard group, usage \Pic(X)
% \Picardstack_{X/B} denotes the Picard stack of X over B
% \Picardfunctor_{X/B} denotes the Picard functor of X over B
\def\Pic{\mathop{\mathrm{Pic}}\nolimits}
\def\Picardstack{\mathcal{P}\!\mathit{ic}}
\def\Picardfunctor{\mathrm{Pic}}
\def\Deformationcategory{\mathcal{D}\!\mathit{ef}}


\begin{document}

\title{History and Philosophy of Biology}
\maketitle

\phantomsection
\label{section-phantom}

\tableofcontents

\section{What this chapter covers}
\label{section-overview}

These notes follow the history of
biological thought as Daniel's first-year
course told it: Aristotle; the early modern
natural philosophers (Bacon, Descartes,
Leibniz, Newton); Linnaeus and Buffon;
Lamarck, Cuvier, Geoffroy and Lyell; Darwin
and Wallace, the struggle for existence and
the debate on whether natural selection is
a true cause; the objection from blending
inheritance; and the road from Mendel to
the modern synthesis and the central dogma.
A last section records the thesis of an
essay Daniel wrote later on evolution as
ordered, directed change.

\textbf{Sources of these notes.}
The notes rest on the course
\emph{Filosof\'ia e Historia de la
Biolog\'ia} (first semester, Facultad de
Ciencias, UNAM, autumn 2013). What survives
is team-written work in which Daniel took
part: answers to three take-home exams
(Aristotle, Bacon, Descartes and Leibniz;
Lamarck, Cuvier, Darwin and extinction; a
long essay on Darwin and the \emph{vera
causa}), a lab report on Linnaeus's sexual
system of plants, film reviews, and the
course slides on Aristotle, Leibniz,
chapter III of the \emph{Origin} and a
chapter on genes and organisms from the
Spanish book \emph{Trayectorias de vida}.
The last section rests on a two-page essay
from the course \emph{Evoluci\'on} (seventh
semester). The facts were checked and
completed with Jaeger's \emph{Aristotle}
\cite{jaeger1948aristotle}, Darwin's
\emph{Origin} in Beer's edition
\cite{darwin2008origin}, Mayr's \emph{What
Evolution Is} \cite{mayr2002evolution},
chapters of the UNAM compilation
\emph{Fundamentos hist\'oricos de la
biolog\'ia}
\cite{ruiz2008metodo, valor2008cartesiana,
sloan2008buffon, bueno2008lyell,
ledesma2008lamarck, zamora2008forma},
chapters of the Spanish anthology
\emph{Filosof\'ia e historia de la
biolog\'ia} read in the course
\cite{sloan-natural-history,
guillaumin-vera-causa, bowler-salida,
mayr-fisicos, gould-jenkin,
bowler-mendelismo, bowler-mendel,
barahona-biometrica, barahona-morgan,
kingsland-neodarwinismo}, Singh's
\emph{Plant Systematics}
\cite{singh2010plant}, Hickman et al.\
\cite{hickman2002animal}, and Levins and
Lewontin \cite{levins1985dialectical}.

\section{Aristotle's biology}
\label{section-aristotle}

The course divided the Aristotelian corpus
into four groups: (I) the \emph{Organon}
(logic); (II) the physical world
(\emph{Physics}, \emph{On the Heavens},
\emph{Meteorology}, \emph{On Generation and
Corruption}); (III) the living world;
(IV) \emph{Metaphysics}, \emph{Ethics},
\emph{Politics}, \emph{Poetics}. Group III
splits into observations, the
\emph{History of Animals}, and theory:
\emph{Parts of Animals}, \emph{Progression
of Animals}, \emph{On the Soul}, the
\emph{Parva Naturalia}, \emph{Movement of
Animals} and \emph{Generation of Animals}.

Jaeger places the zoological works late, in
the period of organized research of
Aristotle's school. The first book of
\emph{Parts of Animals} rejects Plato's
method of division by dichotomy as useless
for zoology, since it tears related species
apart; the new classification was to be
wrested from the animals themselves. For
Jaeger the real achievement was not the
classification but that ``for the first
time the observation and description of the
individual and its life-history was being
taken absolutely seriously''
\cite[Chapter XIII]{jaeger1948aristotle}.

The method and ideas, as the course
presented them:
\begin{itemize}
\item Credit to observation and to
classification, against Plato's truth
reached a priori by recollection.
\item The four causes: material, formal,
efficient and final. A difference of form,
function or behaviour is explained when its
four causes are found; the final cause
(what a part is for) comes first.
\item Three kinds of soul: vegetative
(plants), sensitive (animals) and rational
(humans), each including the ones before.
\item Logic as a tool: classification by
inclusion of classes in syllogisms (``the
dog has hair; animals with hair are
mammals; so the dog is a mammal'').
\item Parts are homoeomerous (uniform, like
flesh, bone, blood) or anhomoeomerous
(organs). Parts are compared ``by more or
less'' and ``by analogy'', forerunners of
homology and analogy.
\item Animals divide into \emph{enaima},
with blood (roughly the vertebrates), and
\emph{anaima}, bloodless (roughly the
invertebrates).
\item The \emph{scala naturae}: a graded
series from the inanimate through plants
and animals to man. Species are fixed;
some animals arise by spontaneous
generation.
\end{itemize}

\section{Bacon, Descartes, Leibniz, Newton}
\label{section-early-modern}

The course treated these four as the
background against which modern biology
took its methods.

\textbf{Bacon} (1561--1626). His
\emph{Novum Organum} (1620) belongs to the
\emph{Great Instauration}. The
four \emph{idols} are the obstacles to
seeing nature as it is:
\begin{itemize}
\item idols of the \emph{tribe}: tendencies
common to all humans (trust in the senses,
simplifying what is complicated);
\item idols of the \emph{cave}: the biases
of each individual;
\item idols of the \emph{marketplace}
(forum): errors from language and
ambiguous words;
\item idols of the \emph{theatre}: errors
from philosophical systems and false rules
of demonstration.
\end{itemize}
His induction goes from particulars to
general laws through three tables:
presence (where the phenomenon occurs),
absence (where it does not) and degrees
(how it varies), helped by ``prerogative
instances''. Knowledge should be useful;
Bacon rejected the Greek and Latin
authorities of the universities. Ruiz and
Ayala note that the idea that science
advances by accumulating facts and
extracting a theory from them goes back to
Bacon, and that it is false
\cite[Chapter 3]{ruiz2008metodo}.

\textbf{Descartes} (1596--1650).
\emph{Mechanism}: natural phenomena can and
must be explained by matter, motion and
their laws alone. In the
\emph{Principles of Philosophy} (1644) the
nature of body is extension alone
\cite[Chapter 5]{valor2008cartesiana}. The
course stressed the clock metaphor: living
beings are machines whose gears are too
small to see; God is the clockmaker. The
metaphor removes souls from the working of
bodies, invites the study of their parts,
and supports fixism (a divine machine has
no imperfect species).

\textbf{Leibniz} (1646--1716). Two
principles, central in the course:
\begin{itemize}
\item \emph{Identity of indiscernibles}: no
two substances differ only in place, time
and number; something intrinsic tells them
apart. Illustration: a princess sends her
servants to find two identical leaves in
her garden, and they find none.
\item \emph{Continuity} (\emph{natura non
facit saltus}): species are so alike that
the senses cannot fix where one ends and
the next begins.
\end{itemize}
For Leibniz a species is defined by common
generation, not by form. A newborn child of
monstrous shape was declared human
only ``provisionally''; what counts is
the inner nature, judged by conjecture,
while appearances are judged directly. He
allowed that related forms might descend
from a common ancestral species. Sloan
describes his dynamic nature, governed by
internal laws, and his attack on Newton's
absolute space and time in the
Leibniz--Clarke correspondence (1717)
\cite{sloan-natural-history}.

\textbf{Newton} (1642--1727). His rules of
reasoning open Book III of the
\emph{Principia}. Rule I: admit no more
causes of natural things than those that
are true and sufficient to explain the
phenomena. A cause without solid evidence
is a \emph{hypothesis}, to be banished
from natural philosophy. Rule I became the
methodology of the \emph{vera causa} (true
cause) that dominated British science in
the first half of the nineteenth century
\cite{guillaumin-vera-causa}.

\section{Linnaeus and Buffon}
\label{section-linnaeus-buffon}

\textbf{Linnaeus} (1707--1778), a Swedish
botanist at Uppsala, ordered species into
genera, genera into orders and orders into
classes in his \emph{Systema Naturae}. The
modern zoological hierarchy has seven
mandatory ranks: kingdom, phylum, class,
order, family, genus, species
\cite[Chapter 4]{hickman2002animal}. He set
up binomial nomenclature, and his
\emph{Species Plantarum} (1753) is the
starting point of botanical nomenclature.
His \emph{sexual system} has 24 classes,
defined mostly by the stamens (number,
relative length, fusion), with orders by
the carpels (Monogynia, Digynia, \dots).
Linnaeus knew it was convenient rather than
natural \cite[Chapter 10]{singh2010plant}.
He was a fixist; the course quoted the
motto ``Deus creavit, Linnaeus disposuit''.

\begin{center}
\begin{tabular}{ll}
Class & Character \\
\hline
1--13 & number of stamens \\
14 Didynamia & 2 long, 2 short \\
15 Tetradynamia & 4 long, 2 short \\
16 Monadelphia & filaments in 1 bundle \\
17--18 & in 2 or more bundles \\
19 Syngenesia & united by anthers \\
20 Gynandria & united to gynoecium \\
21 Monoecia & monoecious plants \\
22 Dioecia & dioecious plants \\
23 Polygamia & polygamous plants \\
24 Cryptogamia & no flowers
\end{tabular}
\end{center}

In the course's lab the students removed
petals and looked at stamens and pistils of
three flowers under the microscope,
then followed Linnaeus's key. One flower
had four stamens, two long and two short
(Didynamia); one had its filaments united
in a single bundle (Monadelphia); the third
plant bore separate male and female flowers
(Monoecia). These are classes of the
system, not families. The report concluded
that the key is quick to use without
training, but needs flowers and stops far
above the species.

\textbf{Buffon} (1707--1788) opened his
\emph{Histoire naturelle} (first volume
1749) with the \emph{Premier discours}, an
attack on all taxonomy of the time and on
Linnaeus in particular. Systems built on
essential characters impose an arbitrary
order of the mind; ``there actually exist
in nature only individuals, and the
Genera, Orders and Classes exist only in
our imagination''; nature descends ``by
almost insensible degrees'', so no perfect
classification is possible; and we know
relations, not essences
\cite{sloan2008buffon}. The course added
that he found the sexual system artificial,
built on a single character, and held that
the naturalist should explain nature, not
only order it.

\section{Lamarck, Cuvier, Geoffroy, Lyell}
\label{section-lamarck-cuvier-lyell}

\textbf{Lamarck} (1744--1829) entered the
Jardin du Roi as a botanist with Buffon's
support. His \emph{Philosophie zoologique}
(1809) gave the first coherent theory of
evolution, with two factors: a natural
tendency of organisms towards complexity,
which explains the graded series of
classifications, and the influence of
circumstances, which explains the
deviations from it. Its two laws:
\begin{enumerate}
\item In an animal still developing,
frequent use of an organ strengthens and
enlarges it; constant disuse weakens it
until it disappears.
\item What individuals gain or lose in this
way is passed to their offspring, provided
the change is common to both parents.
\end{enumerate}
Lamarck coined the word \emph{biology}, and
the term \emph{invertebrates} is due to
him \cite{ledesma2008lamarck}.

\textbf{Cuvier} (1769--1832) proved
extinction: the mammoth is a species
distinct from the living African and
Indian elephants, so it is extinct. He
remained a fixist and explained
extinctions by \emph{catastrophism}:
repeated revolutions of the Earth's
surface that destroyed whole faunas. His
anatomy was \emph{functionalist}: the
structure of an animal is to be explained
by its \emph{conditions of existence}, the
adaptive functions of its parts
\cite{zamora2008forma}.

\textbf{Geoffroy Saint-Hilaire}
(1772--1844), Cuvier's colleague at the
Paris Museum, held the opposite
\emph{formalist} view: a group's anatomy
is to be explained by a basic plan or
archetype (unity of type,
theory of analogues). The two clashed in a
famous public debate at the Paris Academy
of Sciences in 1830
\cite{zamora2008forma}. Darwin later
reconciled them: ``unity of type is
explained by unity of descent'', while the
conditions of existence are ``fully
embraced by the principle of natural
selection'' and form ``the higher law''
\cite[Chapter VI]{darwin2008origin}.

\textbf{Lyell} (1797--1875). His
\emph{Principles of Geology} (1830--1833)
rest on \emph{uniformitarianism} (a word of
Whewell, 1832): the past is explained by
the causes we see acting today. Rudwick
distinguished four senses in Lyell:
uniformity of law; uniformity of process
(actualism); gradualism (slow, constant
change); and uniformity of state, or
non-progressionism. The first two are
methodological; the last two are
hypotheses about the Earth
\cite{bueno2008lyell}. Mountains, for
instance, are not raised by an unknown
force but by causes observed now: gradual
uplift and volcanic action.

\textbf{Extinction.} One thread of the
course's second exam:
\begin{center}
\begin{tabular}{ll}
Author & Extinction \\
\hline
Aristotle & none: fixed chain \\
Linnaeus & none: species created \\
Buffon & species change \\
Cuvier & real, by catastrophes \\
Lamarck & species transform \\
Darwin & natural, by competition
\end{tabular}
\end{center}

\section{Darwin, Wallace, natural selection}
\label{section-darwin-wallace}

Darwin read Malthus's \emph{Essay on the
Principle of Population} in September
1838, near the end of his notebook D.
Wallace reached the same idea in 1858 on
Gilolo (not, as he later said, on Ternate)
and sent his paper to Darwin, who received
it on 12 June 1858. On the advice of Lyell
and Hooker, Wallace's paper and an extract
of Darwin's were read together at the
Linnean Society \cite{bowler-salida}. The
\emph{Origin} appeared in 1859.

\textbf{Natural selection.} Owing to the
struggle for life, any variation, however
slight, if in any degree profitable to an
individual, ``will tend to the preservation
of that individual, and will generally be
inherited by its offspring''. Darwin named
the principle by analogy with man's
selection
\cite[Chapter III]{darwin2008origin}.

\textbf{The struggle for existence}
(\emph{Origin}, chapter III, summarized
in the course's slides):
\begin{itemize}
\item The term is used ``in a large and
metaphorical sense'', including dependence
of one being on another and success in
leaving progeny: a plant at the edge of a
desert struggles against drought.
\item It follows from geometrical increase:
more individuals are born than can
survive, so they struggle with their own
species, with other species, or with
physical conditions. It is ``the doctrine
of Malthus applied with manifold force''.
\item Linnaeus's annual plant with only two
seeds would give a million plants in
twenty years. The elephant, slowest
breeder, breeding from 30 to 90 years old
and bearing three pairs of young, would
leave fifteen million descendants of one
pair after five centuries.
\item Checks: food, predators, climate,
epidemics. The number of eggs matters less
than the destruction at some stage of
life.
\item The struggle is most severe between
individuals and varieties of the same
species, and often severe between species
of the same genus, which share structure
and habits.
\end{itemize}
\cite[Chapter III]{darwin2008origin}.
The plant computation is
$$
2^{20} = 1\,048\,576 \approx 10^{6}.
$$

\section{Natural selection as a true cause}
\label{section-vera-causa}

This was the theme of the course's main
essay, written by Daniel's team, on the
methodology of the \emph{vera causa}.

\textbf{Herschel} (1792--1871), in his
\emph{Preliminary Discourse}, faced causes
that cannot be observed (the nature of
light, the slow raising of mountains). To
accept such a cause one must show that it
really exists, that it is sufficient to
produce the phenomenon, and that there is
independent evidence for it; a surprising
prediction counts strongly in its favour.
\textbf{Whewell} (1794--1866) gave
hypotheses a real role in discovery and
accepted them for unexpected predictions,
simplicity and the \emph{consilience of
inductions}: explaining different empirical
generalizations at once, as gravitation
explains both Kepler's laws and Galileo's
kinematics \cite{guillaumin-vera-causa}.
Whewell also separated historical from
mechanical causation.

Darwin followed Herschel. He argued for
natural selection by analogy with breeders:
artificial selection is a real cause, and
natural selection acts in the same way.
But the analogy shows the reality of the
cause, not its sufficiency: breeders make
races, not species, in short times. This
is the point Fleeming Jenkin attacked in
1867 \cite{guillaumin-vera-causa}. The
team's essay concluded that natural
selection failed Newton's test (it was not
observed), was questioned under Herschel's
(sufficiency) and fell between Whewell's
two kinds of causation; yet, as with
Lyell's mountains, the successful theory
was kept and the methodology changed. The
\emph{Origin} opens with an epigraph from
Whewell: events are brought about ``not by
insulated interpositions of Divine power,
exerted in each particular case, but by
the establishment of general laws''
\cite{darwin2008origin}.

\textbf{Teleology and mechanism.} Final
causes passed from Aristotle to natural
theology, where God was the final cause of
everything. Boyle (1688) saw that
explaining how a structure works
mechanically is a different task from
explaining why it exists: \emph{proximate}
versus \emph{ultimate} causes. Mayr's
example: the proximate cause of sexual
dimorphism in plumage is hormonal, the
ultimate cause is sexual selection
\cite{mayr-fisicos}. Darwin moved from a
designed nature to one without design:
natural selection improves lineages but
aims at no final perfect form.

\section{Blending inheritance and Jenkin}
\label{section-blending}

Darwin and his critics shared the usual
theory of \emph{blending inheritance}:
offspring are intermediate between their
parents. Fleeming Jenkin, an engineer,
argued in the \emph{North British Review}
(1867) that a single favourable variant
would be swamped by crossing with the
common form, so it could never spread
(his example, a shipwrecked white man on
an island, is a reminder of Victorian
racism). Gould shows that Darwin had long
known the problem. What Jenkin did was to
separate two kinds of variation: blending
swamps rare single large variations
(``sports''), but not the small,
recurrent individual differences. Darwin
welcomed the argument, which strengthened
his preference for small continuous
variation. Jenkin, a close friend of
Kelvin, also pressed Kelvin's young age of
the Earth \cite{gould-jenkin}. Particulate
(Mendelian) inheritance removes the
difficulty, since genes do not dilute; in
1930 Fisher dismissed blending as
irrelevant \cite{kingsland-neodarwinismo}.

\section{After Darwin: heredity, synthesis}
\label{section-after-darwin}

\textbf{Mendel.} His hybridization
experiments with peas, grown in the
monastery garden at Brno, were presented
to the Brno Natural History Society and
published in its journal in 1866. They
drew little attention; Mendel's only
correspondence with a major biologist was
with N\"ageli, from 1866 to 1873
\cite{bowler-mendel}.

\textbf{Rediscovery} (1900) by Correns, de
Vries and Tschermak; Bateson made
Mendelism the basis of a new science of
heredity. Whether de Vries and Tschermak
understood Mendel's laws independently is
disputed; Correns is the least
controversial of the three
\cite{bowler-mendelismo}.

\textbf{Biometricians against Mendelians.}
Galton (continuous variation, regression
towards the ancestral mean), Pearson and
Weldon defended continuous evolution and
ancestral heredity; they founded
\emph{Biometrika} in 1901. Bateson held
that evolution proceeds by discontinuous
variation \cite{barahona-biometrica}. The
early Mendelians (de Vries, Bateson,
Johannsen) were saltationists
\cite[Chapter 4]{mayr2002evolution}.

\textbf{Morgan.} Sutton had noted in 1902
the parallel between chromosomes and
Mendelian factors. In 1910 Morgan found a
white-eyed male \emph{Drosophila
melanogaster} and explained its pattern by
genes on the X chromosome (sex-linked
inheritance). With Sturtevant, Muller and
Bridges he wrote \emph{The Mechanism of
Mendelian Heredity} (1915), which set out
the chromosome theory of heredity; Morgan
received the Nobel Prize in 1933
\cite{barahona-morgan}.

\textbf{The modern synthesis.} Fisher's
\emph{Genetical Theory of Natural
Selection} (1930) joined Mendelism and
Darwinism; with Haldane and Wright he
founded population genetics.
Dobzhansky's \emph{Genetics and the Origin
of Species} (1937) was followed by Mayr
(1942), Simpson (1944) and Stebbins
(1950); Huxley's book of 1942 named the
``modern synthesis''
\cite{kingsland-neodarwinismo}. Mayr calls
the acceptance of gradualism, speciation
and natural selection in the 1940s the
second Darwinian revolution, the first
being the quick acceptance of evolution
and common descent after 1859
\cite[Chapter 4]{mayr2002evolution}.

\textbf{The central dogma.} The course's
slides on \emph{Trayectorias de vida}
went from Mendel and Galton through
Morgan to Watson and Crick, Gamow's
codons, the central dogma (DNA makes RNA,
RNA makes protein) and Jacob and Monod's
operons, and warned against ``genes for''
traits and biological determinism. Mayr
reads the dogma as the final refutation of
soft inheritance: no information passes
from the proteins of the body to the
nucleic acids of the germ cells
\cite[Chapter 4]{mayr2002evolution}.

\section{Essay: evolution as directed change}
\label{section-essay}

In the course \emph{Evoluci\'on} Daniel
wrote a short response to a chapter
titled ``Evolution as Theory and
Ideology'', the first chapter of Levins
and Lewontin's \emph{The Dialectical
Biologist}. There evolutionism is a world
view built on the hierarchy of concepts
change, order, direction, progress and
perfectability
\cite[Chapter 1]{levins1985dialectical}.

The essay's thesis:
\begin{itemize}
\item Evolution should first be defined as
a general phenomenon, not only a
biological one, with three properties:
change, order and direction.
\item \emph{Order}: differences between
elements matter only if they place them at
different points of some scale.
\item \emph{Direction} makes evolution
unidirectional rather than cyclic.
Complexity is rejected as the direction,
since it flatters our own place.
\item \emph{Stability} is a better
candidate: evolution as a property of a
dynamic system that may be in a steady
state or not; a dynamic equilibrium, whose
parts never stop moving, can look unstable
from a reductionist point of view.
\item Darwin's materialism against Plato's
essentialism: small differences between
individuals become differences between
species over long times, so species and
individuals are related ontologically;
Plato's species differ by essence.
\end{itemize}
The essay links back to the course's
Leibniz and Buffon: continuity of forms
and the reality of individuals.

\bibliography{my}
\bibliographystyle{amsalpha}

\end{document}
